% Faithful formal-publication transcription of Sections4--5, including all local continuations.
\section*{4 Proof of Leaf Lemma}
We prove Leaf Lemma in three subsections.
(Section4.1) We start by recalling a distributional characterisation due to Blais and Brody [7] of expected query complexity $\overline R_\epsilon$ using decision trees that can ``abort''.
(Section4.2) We then formulate a Hard Side Lemma, which encapsulates the core challenge in finding the hard side of a given function $g$ and from which Leaf Lemma is easy to derive.
(Section4.3) Finally, we prove the Hard Side Lemma.
\subsection*{4.1 Distributional characterisation of $\overline R_\epsilon$ due to Blais--Brody}
A (deterministic) abort-tree $t$ is a decision tree that outputs either a boolean value (0 or 1) or the abort symbol $\bot$. When an abort-tree is trying to compute a boolean function $g$, we do not consider the output $\bot$ as an ``error''; the tree simply gives up on the computation. Indeed, we say that $t(x)$ errs iff $t(x)=1-g(x)$, that is, $t(x)\ne\bot$ and $t(x)\ne g(x)$. As before, a randomised abort-tree is a probability distribution over deterministic abort-trees. For $\gamma\in(0,1)$ and $\epsilon\in[0,1/2)$ we define $R_{\gamma,\epsilon}(g)$ as the least (worst-case) depth of a randomised abort-tree $T$ such that for all $x\in\operatorname{dom}(g)$:
\[\Pr_{t\sim T}[t(x)=\bot]\le\gamma\quad\text{and}\quad\Pr_{t\sim T}[t(x)\text{ errs}]\le\epsilon.\]
We formulate a distributional version of $R_{\gamma,\epsilon}(g)$ as follows. For a distribution $\mu$ over $\operatorname{dom}(g)$, we define $D^\mu_{\gamma,\epsilon}(g)$ as the least depth of a deterministic abort-tree $t$ such that
\[\Pr_{x\sim\mu}[t(x)=\bot]\le\gamma\quad\text{and}\quad\Pr_{x\sim\mu}[t(x)\text{ errs}]\le\epsilon.\]
The following two lemmas from [7, Section3.1] connect abort-trees and $\overline R_\epsilon(g)$.
\begin{authorclaim}{Lemma8 (Abort vs. expected depth)}
For every $\epsilon\in[0,1/2)$ and $\gamma\in(0,1)$,
\[\gamma\cdot R_{\gamma,\epsilon}(g)\le\overline R_\epsilon(g)\le\frac1{1-\gamma}\cdot R_{\gamma,(1-\gamma)\epsilon}(g).\]
\end{authorclaim}
\begin{authorclaim}{Lemma9 (Minimax)}
For every $\epsilon\in[0,1/2)$, $\gamma\in(0,1)$, and $\alpha,\beta\in(0,1)$ with $\alpha+\beta\le1$,
\[\max_\mu D^\mu_{\gamma/\alpha,\epsilon/\beta}(g)\le R_{\gamma,\epsilon}(g)\le\max_\mu D^\mu_{\alpha\gamma,\beta\epsilon}(g).\]
\end{authorclaim}
\subsection*{4.2 Statement of Hard Side Lemma}
When searching for the hard side of a partial function $g$ under a distribution $\mu=\frac12(\mu^0+\mu^1)$, it is convenient to study a more symmetric notion of noisiness than the one-sided variant defined earlier. For a leaf $\ell\in L(t)$ of an abort-tree $t$, we define the relative error $\operatorname{re}(\ell,\mu)$ so that if $\ell$ is an aborting leaf, then $\operatorname{re}(\ell,\mu):=0$; otherwise
\[\operatorname{re}(\ell,\mu):=\frac{\min\{\ell(\mu^0),\ell(\mu^1)\}}{\ell(\mu^0)+\ell(\mu^1)}\in[0,1/2].\]
This definition captures the best achievable error of a leaf in an abort-tree. Namely, let us say that $t$ is $\mu$-smart if every non-abort leaf $\ell\in L(t)$ outputs a boolean value $i\in\{0,1\}$ that maximises $\ell(\mu^i)$. Then for every leaf $\ell$ in a $\mu$-smart $t$ we have $\Pr_{x\sim\mu}[t(x)\text{ errs}\mid\ell_x^t=\ell]=\operatorname{re}(\ell,\mu)$. An easy calculation gives the following claim, which we record for future use.
\begin{authorclaim}{Claim10}
For a $\mu$-smart $t$ we have $\Pr_{x\sim\mu}[t(x)\text{ errs}]=\mathbb E_{x\sim\mu}[\operatorname{re}(\ell_x^t,\mu)]$.
\end{authorclaim}
Another easy calculation shows that relative error implies noisiness.
\begin{authorclaim}{Claim11}
If $\operatorname{re}(\ell,\mu)\ge\epsilon$, then $\ell$ is $(\epsilon,\mu,b)$-noisy for both $b\in\{0,1\}$.
\end{authorclaim}
We are now ready to formulate Hard Side Lemma, which isolates the technical challenge 3.1 (discussed in Subsection3.1): Every partial function $g$ admits a balanced distribution $\mu$ and a hard side $b$ such that if we run a shallow abort-tree on the hard side $\mu^b$ of $\mu$, then $t$ must either abort with high probability or we reach a leaf of noticeable error (in expectation).
\begin{authorclaim}{Hard Side Lemma}
For every partial function $g$ and $0<\epsilon\lll\delta\ll1$, there exists a distribution $\mu=\frac12(\mu^0+\mu^1)$ over $\operatorname{dom}(g)$ and a hard side $b\in\{0,1\}$ such that for any deterministic abort-tree $t$ with $\operatorname{depth}(t)/\overline R_\epsilon(g)\lll\delta$ we have either
\begin{equation}\tag{3}\Pr_{x\sim\mu^b}[t(x)=\bot]>1-\delta\quad\text{or}\quad\mathbb E_{x\sim\mu^b}[\operatorname{re}(\ell_x^t,\mu)]>\epsilon.\end{equation}
\end{authorclaim}
We defer the proof until Subsection4.3. We first use the lemma to prove Leaf Lemma, and here is where we address challenge 3.1: we exploit the high abort probability (namely, $1-\delta$) guaranteed by Hard Side Lemma to show that typical leaves are noisy.
\begin{authorclaim}{Leaf Lemma (restated)}
For every partial $g$ and $0<\epsilon\lll\delta\ll1$, there exists a distribution $\mu=\frac12(\mu^0+\mu^1)$ over $\operatorname{dom}(g)$ and a hard side $b\in\{0,1\}$ such that for every deterministic tree $t$ and $i\in\{0,1\}$:
\[\frac{\mathbb E_{x\sim\mu^i}[q(t,x)]}{\overline R_\epsilon(g)}\lll\delta\quad\Longrightarrow\quad\Pr_{x\sim\mu^i}[\ell_x^t\text{ is }(\epsilon,\mu,b)\text{-noisy}]\ge1-\delta.\]
\end{authorclaim}
\begin{proof}
We observe first that regardless of $\mu,b$, or even the expected depth of $t$, the lemma holds for the easy side $i=1-b$. Indeed, if we let $B\subseteq L(t)$ denote the set of non-$(\epsilon,\mu,b)$-noisy leaves,
\[\Pr_{x\sim\mu^{1-b}}[\ell_x^t\in B]=\sum_{\ell\in B}\ell(\mu^{1-b})<\epsilon\sum_{\ell\in B}\ell(\mu^b)\le\epsilon\sum_{\ell\in L(t)}\ell(\mu^b)=\epsilon\le\delta.\]
Let us then focus on the interesting case $i=b$ where the careful choice of $\mu$ and $b$ is essential. We invoke Hard Side Lemma with parameters $\epsilon$ and $\dot\delta:=\delta^2$ (assuming suitably $0<\epsilon\lll\dot\delta\ll1$) to obtain $\mu$ and $b$ such that for every abort-tree $\dot t$ with $\operatorname{depth}(\dot t)\le\dot\delta^c\overline R_\epsilon(g)$ the property(3) holds (with dotted parameters). Let $x\sim\mu^b$ henceforth and write $\operatorname{re}(\ell):=\operatorname{re}(\ell,\mu)$ for short. Suppose $t$ satisfies $\mathbb E_x[q(t,x)]\le\dot\delta^{c+2}\overline R_\epsilon(g)$ (where we chose $2(c+2)$ as the exponent hidden by $\lll$). Recalling from Claim11 that relative error implies noisiness, our goal is to show
\begin{equation}\tag{4}\Pr_x[\operatorname{re}(\ell_x^t)\ge\epsilon]\ge1-\delta.\end{equation}
We convert $t$ into an abort-tree by letting $t'$ be a modification of $t$ that aborts whenever more than $\dot\delta^c\overline R_\epsilon(g)$ queries are made. Using Markov's inequality and the low expected depth of $t$,
\[\Pr_x[t'(x)=\bot]=\Pr_x[q(t,x)>\dot\delta^c\overline R_\epsilon(g)]\le\frac{\mathbb E_x[q(t,x)]}{\dot\delta^c\overline R_\epsilon(g)}\le\dot\delta^2.\]
We also have $\Pr_x[\operatorname{re}(\ell_x^t)\ge\epsilon]\ge\Pr_x[\operatorname{re}(\ell_x^{t'})\ge\epsilon]$ since we only made more executions abort. To prove(4), suppose for contradiction that $\Pr_x[\operatorname{re}(\ell_x^{t'})\ge\epsilon]<1-\delta$. Let $\dot t$ be a further modification of $t'$ that aborts any leaf $\ell\in L(t')$ with $\operatorname{re}(\ell)\ge\epsilon$. Note that
\[\Pr_x[\dot t(x)=\bot]\le\Pr_x[t'(x)=\bot]+\Pr_x[\operatorname{re}(\ell_x^{t'})\ge\epsilon]\le\dot\delta^2+1-\delta\le1-\dot\delta.\]
Hence we get from (dotted) property(3) that $\mathbb E_x[\operatorname{re}(\ell_x^{\dot t})]>\epsilon$. But this contradicts the fact that $\operatorname{re}(\ell)<\epsilon$ for all $\ell\in L(\dot t)$ by construction. This verifies(4) and concludes the proof.
\end{proof}
\subsection*{4.3 Proof of Hard Side Lemma}
Let $\nu$ be a distribution that witnesses $D:=\max_{\nu'}D^{\nu'}_{1-\delta,\epsilon^{1/3}}(g)$ so that every abort-tree $t$ with $\operatorname{depth}(t)<D$ fails to satisfy at least one of the following:
\begin{align}\Pr_{x\sim\nu}[t(x)=\bot]&\le1-\delta,\tag{5}\\\Pr_{x\sim\nu}[t(x)\text{ errs}]&\le\epsilon^{1/3}.\tag{6}\end{align}
As a minor technicality, we re-balance $\nu$. We can write $\nu=\lambda\mu^0+(1-\lambda)\mu^1$ where $\lambda\in(0,1)$ and $\mu^i$ is a distribution supported on $g^{-1}(i)$. We define $\mu:=\frac12(\mu^0+\mu^1)$ as our balanced distribution.
Assume towards a contradiction that there does not exist a hard side for $\mu$, that is, the claim of the lemma fails for both $b\in\{0,1\}$. This means there exists two abort-trees $t_0$ and $t_1$ of depth at most $\delta^3\overline R_\epsilon(g)$ (where we chose 3 as the exponent hidden by $\lll$) such that for both $b\in\{0,1\}$:
\begin{align}\Pr_{x\sim\mu^b}[t_b(x)=\bot]&\le1-\delta,\tag{7}\\\mathbb E_{x\sim\mu^b}[\operatorname{re}(\ell_x^{t_b},\mu)]&\le\epsilon.\tag{8}\end{align}
We will use $t_0$ and $t_1$ to construct a third tree $t$ that computes $g$ too well relative to $\nu$ contradicting our choice of $\nu$. We may assume wlog that $t_0$ and $t_1$ are $\mu$-smart, since the properties(7)--(8) do not depend on the boolean leaf-labels (only whether a leaf aborts or not). We now define $t$ as follows: On input $x$ we run both $t_0(x)$ and $t_1(x)$; if $t_0(x)\ne\bot$, we output $t_0(x)$; otherwise we output $t_1(x)$. We will show that $t$ has $\operatorname{depth}(t)<D$ and satisfies(5)--(6), which will contradict our choice of $\nu$.
\paragraph{Tree $t$ is shallow.} We have the following chain of inequalities
\begin{align*}
\operatorname{depth}(t)&\le2\delta^3\overline R_\epsilon(g)\le32\delta^3\overline R_{\epsilon^{1/4}}(g)<\delta^2\overline R_{\epsilon^{1/4}}(g)\le R_{1-\delta^2,\delta^2\epsilon^{1/4}}(g)\\
&\le\max_{\nu'}D^{\nu'}_{(1-\delta^2)^2,\delta^4\epsilon^{1/4}}(g)\le\max_{\nu'}D^{\nu'}_{1-\delta,\epsilon^{1/3}}(g)=:D.
\end{align*}
The first inequality uses the definition of $t$. Second uses error reduction (Claim7 with $k:=4$). Third uses $\delta\ll1$. Fourth uses Lemma8 (with $\gamma:=1-\delta^2$). Fifth uses the minimax lemma (Lemma9 with $\alpha:=1-\delta^2$, $\beta:=\delta^2$). The final inequality uses $\epsilon\lll\delta\ll1$.
\paragraph{Tree $t$ has bounded abort.} We verify property(5) by
\begin{align*}
\Pr_{x\sim\nu}[t(x)=\bot]&=\Pr_{x\sim\nu}[t_0(x)=\bot\wedge t_1(x)=\bot]\\
&=\lambda\Pr_{x\sim\mu^0}[t_0(x)=\bot\wedge t_1(x)=\bot]+(1-\lambda)\Pr_{x\sim\mu^1}[t_0(x)=\bot\wedge t_1(x)=\bot]\\
&\le\lambda\Pr_{x\sim\mu^0}[t_0(x)=\bot]+(1-\lambda)\Pr_{x\sim\mu^1}[t_1(x)=\bot]\le1-\delta.\quad\text{(Using(7))}
\end{align*}
\begin{center}\includegraphics[width=.9\linewidth]{sources/author-figure1.pdf}\end{center}
Figure1. Two trees $t_0$ and $t_1$ in the proof of Hard Side Lemma. The leaves partition $\operatorname{dom}(g)$ into subcubes where grey leaves output $\bot$, green leaves output 1, and blue leaves output 0. Hatched regions are error. We are promised that, e.g., $t_0$ has bounded abort (10\% in our figure) over $\mu^0$, but not necessarily over $\mu^1$.
\paragraph{Tree $t$ errs rarely.} We start with a claim that says that if the expected relative error is low over one side $\mu^b$ of $\mu$, then a $\mu$-smart tree errs rarely over the whole distribution $\mu$.
\begin{authorclaim}{Claim12}
Let $t'$ be $\mu$-smart and $b\in\{0,1\}$. If $\mathbb E_{x\sim\mu^b}[\operatorname{re}(\ell_x^{t'},\mu)]\le\epsilon$ then $\Pr_{x\sim\mu}[t'(x)\text{ errs}]\le\epsilon^{1/2}$.
\end{authorclaim}
\begin{proof}
We prove the claim for $b=0$ as the other case is analogous. Since $t'$ and $\mu$ are fixed, we drop them from notation writing $\operatorname{re}(\ell):=\operatorname{re}(\ell,\mu)$, $\ell_x:=\ell_x^{t'}$, $L:=L(t')$. We argue that relative error on one side of the distribution must spill over to the other side:
\begin{align*}
\mathbb E_{x\sim\mu^0}[\operatorname{re}(\ell_x)]&=\sum_{\ell\in L}\ell(\mu^0)\operatorname{re}(\ell)\ge\sum_{\ell\in L}\ell(\mu^1)\operatorname{re}(\ell)^2\\
&=\mathbb E_{x\sim\mu^1}[\operatorname{re}(\ell_x)^2]\ge\mathbb E_{x\sim\mu^1}[\operatorname{re}(\ell_x)]^2.
\end{align*}
Here the first inequality used $\ell(\mu^0)\ge\ell(\mu^1)\operatorname{re}(\ell)$ (from Claim11) and the second inequality used Jensen's inequality. It follows that $\mathbb E_{x\sim\mu^1}[\operatorname{re}(\ell_x)]\le\mathbb E_{x\sim\mu^0}[\operatorname{re}(\ell_x)]^{1/2}\le\epsilon^{1/2}$ and therefore $\mathbb E_{x\sim\mu}[\operatorname{re}(\ell_x)]\le\epsilon^{1/2}$. The claim then follows from Claim10.
\end{proof}
We now verify property(6), which concludes the proof of Hard Side Lemma.
\begin{align*}
\Pr_{x\sim\nu}[t(x)\text{ errs}]&\le\Pr_{x\sim\nu}[t_0(x)\text{ errs}\vee t_1(x)\text{ errs}]\\
&\le\sum_{b\in\{0,1\}}\Pr_{x\sim\nu}[t_b(x)\text{ errs}]\\
&=\sum_{b\in\{0,1\}}\lambda\Pr_{x\sim\mu^0}[t_b(x)\text{ errs}]+(1-\lambda)\Pr_{x\sim\mu^1}[t_b(x)\text{ errs}]\\
&\le\sum_{b\in\{0,1\}}2\Pr_{x\sim\mu}[t_b(x)\text{ errs}]\le\sum_{b\in\{0,1\}}2\cdot\epsilon^{1/2}\quad\text{(Claim12 and(8))}\\
&=4\epsilon^{1/2}\le\epsilon^{1/3}.\quad(\epsilon\ll1)
\end{align*}
